\documentclass[11pt]{article}
\usepackage[margin=1in]{geometry}
\usepackage{amsmath,amssymb,amsthm,mathtools}
\usepackage{hyperref}
\usepackage{enumitem}
\usepackage[T1]{fontenc}

\newtheorem{theorem}{Theorem}[section]
\newtheorem{proposition}[theorem]{Proposition}
\newtheorem{corollary}[theorem]{Corollary}
\newtheorem{conjecture}[theorem]{Conjecture}
\theoremstyle{definition}
\newtheorem{definition}[theorem]{Definition}
\newtheorem{remark}[theorem]{Remark}

\title{Barrier-Aware Research on \(P\) versus \(NP\):\\
A Formal Audit of a Homological Claim and a\\
Verifier-Invariant Obstruction Program}
\author{ChatGPT}
\date{September 2026}

\begin{document}
\maketitle

\begin{abstract}
We establish an auditable baseline for a research program on \(P\) versus
\(NP\). We do not claim a solution of the problem. Our first contribution is
a formal audit of the proposed homological separation in arXiv:2510.17829.
Four explicit counterexamples invalidate indispensable steps: the proposed
boundary is not closed on directed computation paths; the stated category of
problems and many-one reductions is not additive; a verification-order
average is not additive under composition; and the displayed contracting
homotopy fails on a one-edge deterministic computation. These results are
formalized in Lean.

Our second contribution is a presentation-invariance principle. Adding an
ignored nonempty certificate factor preserves the recognized language while
it may alter raw trace multiplicity. This motivates Verifier-Invariant
Obstruction Theory (VIOT), in which topological data are formed only after
localizing verifier presentations by resource-preserving simulations,
padding, stuttering, and recoding. We formulate a sound obstruction schema
and a barrier-aware program connecting restriction sites, constructive
refuters, convergent gate elimination, proof complexity, and equivariant
topology. The central circuit lower bound remains open.
\end{abstract}

\section{Introduction}

The \(P\)-versus-\(NP\) problem asks whether polynomial-time verification
always entails polynomial-time deterministic decision. By Cook--Levin, a
polynomial-time algorithm for SAT proves \(P=NP\), whereas proving SAT outside
\(P\) proves \(P\ne NP\). Since \(P\subseteq P/poly\), the stronger statement
\[
  \mathrm{SAT}\notin P/poly
\]
also implies \(P\ne NP\).

A credible route must survive local proof checking and global
complexity-theoretic barriers. Relativization \cite{BGS}, natural proofs
\cite{RR}, and algebrization \cite{AW} rule out broad classes of otherwise
persuasive arguments unless an additional ingredient is identified. Explicit
lower bounds for unrestricted fan-in-two Boolean circuits remain linear.
Accordingly, the present paper first falsifies a precise invalid construction
and then states a repaired topological program with every missing bridge made
explicit.

\paragraph{Epistemic status.}
The counterexamples and padding theorem below are proved and formalized. The
VIOT constructions are definitions and conjectures, not a proof of
\(P\ne NP\).

\section{Circuit-lower-bound setting}

A verifier presentation is a tuple \(\mathcal V=(W,V,t)\), where
\(V(x,w)\in\{0,1\}\) is computable within the stated bound and
\[
  L(\mathcal V)=\{x:\exists w,\;V(x,w)=1\}.
\]
For NP, the running time and witness length are polynomially bounded.

Counting shows that almost every Boolean function on \(n\) bits needs
circuits of size \(\Theta(2^n/n)\), but it does not uniformly place a hard
family in NP. Recent range-avoidance work gives near-maximum lower bounds in
exponential-time classes with promise Merlin--Arthur access
\cite{RenWilliams}; descending to an explicit NP family is a central missing
step.

For unrestricted binary circuits, gate elimination remains the main explicit
method. Li and Yang proved a \(3.1n-o(n)\) lower bound over the complete binary
basis. Carmosino, Dang, and Jackman show that simplification is confluent over
DeMorgan and \(\{\land,\lor,\oplus\}\) bases but not over \(\mathcal U_2\)
and \(\mathcal B_2\), and connect simplification with constructive refuters
\cite{CDJ}. Improving a linear constant is valuable but asymptotically
insufficient for NP versus \(P/poly\).

\section{Audit of a claimed homological separation}
\label{sec:audit}

The construction in \cite{Tang} takes directed computation paths
\((c_0,\ldots,c_n)\) as generators and proposes
\[
 d_n(c_0,\ldots,c_n)=
 \sum_{i=0}^n(-1)^i(c_0,\ldots,\widehat c_i,\ldots,c_n).
\]
It also asserts an additive category of problems and reductions, a
verification-order cocycle, and a contraction of deterministic computations.

\subsection{Face closure}

\begin{theorem}[The proposed boundary is not well-defined]
There is a valid directed computation path whose middle face is not a valid
directed path. Hence the displayed \(d_n\) need not map the stated group
\(C_n\) to \(C_{n-1}\).
\end{theorem}

\begin{proof}
Take exactly the transitions
\[
 c_0\longrightarrow c_1\longrightarrow c_2
\]
and no transition \(c_0\to c_2\). Designate \(c_0\) initial and \(c_2\)
terminal. Then \((c_0,c_1,c_2)\) is valid, but deleting \(c_1\) gives
\((c_0,c_2)\), which is not. Thus one summand of
\(d_2(c_0,c_1,c_2)\) is absent from the stated codomain.
\end{proof}

\begin{remark}
The identity \(d^2=0\) cannot repair a typing failure. Directed path homology
uses a face-stable submodule or modified boundary; that would be a different
construction and would still require a complexity bridge.
\end{remark}

\subsection{Additivity}

\begin{theorem}[The claimed zero object does not exist]
With polynomial-time many-one reductions as morphisms, there is no morphism
from a universal nonempty language to the empty language. Consequently the
empty language is not a zero object, and the category is not additive as
claimed.
\end{theorem}

\begin{proof}
A reduction \(f:\Sigma^*\to\varnothing\) would require
\[
 x\in\Sigma^*\quad\Longleftrightarrow\quad f(x)\in\varnothing
\]
for every \(x\). The left side is true and the right side false. Hence the
hom-set is empty. Every hom-set of a preadditive category is an abelian group
and therefore contains a zero morphism.
\end{proof}

The instruction to run two reductions in parallel and combine their outputs
does not define a binary operation, identity, inverse, or compatibility with
composition.

\subsection{Verification-order functional}

For \(\sigma=(\sigma_1,\ldots,\sigma_m)\), \cite{Tang} defines
\[
 \rho(\sigma)=\frac{1}{m-1}
 \sum_{k=1}^{m-1}\operatorname{sgn}(\sigma_{k+1}-\sigma_k)
\]
and uses additivity under composition.

\begin{proposition}
The displayed \(\rho\) is not additive under concatenation.
\end{proposition}

\begin{proof}
For \(a=(1,2)\) and \(b=(3,4)\), one has
\(\rho(a)=\rho(b)=1\). Every adjacent difference in
\(a\|b=(1,2,3,4)\) is positive, so \(\rho(a\|b)=1\), whereas
\(\rho(a)+\rho(b)=2\).
\end{proof}

The cross-boundary term and changed normalization invalidate the claimed
boundary-annihilation calculation.

\subsection{Contracting homotopy}

\begin{proposition}
The proposed append-the-next-configuration homotopy fails on the unique edge
of a one-step deterministic computation.
\end{proposition}

\begin{proof}
Let \(e=[c_0,c_1]\), where \(c_1\) is terminal and
\(d(e)=c_1-c_0\). The rule gives
\(s_0(c_0)=e\), \(s_0(c_1)=0\), and \(s_1(e)=0\). Therefore
\[
 (ds+sd)(e)=s_0(c_1)-s_0(c_0)=-e\ne e.
\]
\end{proof}

\begin{corollary}
The argument in \cite{Tang} does not establish \(P\ne NP\).
\end{corollary}

\begin{proof}
Its proposed boundary is not a map on the stated generators. Independently,
the additive-category, cocycle, and contraction claims fail by the preceding
counterexamples, while the main theorem depends on them.
\end{proof}

\section{Verifier padding and a no-go principle}
\label{sec:padding}

\begin{definition}
For a verifier \(V:X\times W\to\{0,1\}\) and a nonempty set \(D\), define
\[
 V_D(x,(d,w))=V(x,w).
\]
\end{definition}

\begin{theorem}[Padding invariance]
For every \(x\),
\[
 x\in L(V_D)\quad\Longleftrightarrow\quad x\in L(V).
\]
\end{theorem}

\begin{proof}
An accepting padded certificate contains an accepting original certificate.
Conversely, given an accepting \(w\), choose any \(d\in D\).
\end{proof}

\begin{corollary}[Raw-trace no-go principle]
A quantity that changes under ignored nonempty certificate padding does not
factor through the recognized language.
\end{corollary}

Polynomial-time presentation changes can also insert dummy steps, independent
check orders, and finite routing gadgets. Thus trace homology is not an
intrinsic hardness measure until invariance under such moves is proved.

\section{Verifier-Invariant Obstruction Theory}
\label{sec:viot}

\subsection{Localization}

Let \(\mathsf{Ver}\) contain verifier presentations and resource-bounded
simulations. Let \(S\) be generated by certificate recoding, ignored nonempty
padding, polynomial stuttering, and polynomial-time bidirectional simulation.
VIOT seeks invariants on
\[
 \mathsf{Ver}[S^{-1}]
\]
or a homotopy-coherent analogue.

\begin{conjecture}[Simulation localization]
There is an effective localization retaining polynomial resource control and
through which the recognized-language functor factors.
\end{conjecture}

\subsection{Restriction site and defects}

For \(f:\{0,1\}^n\to\{0,1\}\), let \(\mathcal R_n\) be the poset of partial
assignments. For a candidate circuit \(C\), local data over
\(\rho\in\mathcal R_n\) may consist of counterexamples extending \(\rho\),
constructive refuters, simplification certificates, or proof-complexity
witnesses. Restriction maps should form a presheaf \(\mathcal D_{f,C}\).
One then seeks
\[
 \Omega_k(f,C)\in H^k(\mathcal R_n;\mathcal D_{f,C})
\]
or an equivariant index. Recent sign-complex work demonstrates the required
pattern for sign-rank: define a canonical \(\mathbb Z_2\)-space, prove that
its index lower-bounds a standard complexity measure, and compute it for an
explicit family \cite{FHV}. VIOT still needs the analogous circuit-size
bridge.

\begin{theorem}[Abstract obstruction soundness]
Let \(\Omega(f,C)\) take values in a pointed set. Suppose
\[
 C\text{ computes }f\quad\Longrightarrow\quad\Omega(f,C)=0.
\]
If \(\Omega(f,C)\ne0\) for every \(C\) in a class \(\mathcal C\), then no
member of \(\mathcal C\) computes \(f\).
\end{theorem}

\begin{proof}
A circuit in \(\mathcal C\) computing \(f\) would force the obstruction to
be both zero and nonzero.
\end{proof}

This elementary theorem fixes the direction of the bridge: finding cycles is
irrelevant unless exact computation forces vanishing and the target forces
uniform nonvanishing over every candidate.

\subsection{Higher rewriting}

Where circuit simplification is confluent, a normal form may define local
sections. Where confluence fails, selecting a normal form destroys
canonicity. We therefore propose a rewriting 2-groupoid: circuits are objects,
simplifications are 1-cells, and critical-pair resolutions are 2-cells.

\begin{conjecture}[Critical-pair obstruction]
For unrestricted binary circuits, a presentation-invariant equivariant index
of the simplification 2-groupoid lower-bounds a nontrivial circuit resource.
\end{conjecture}

No circuit-size bridge is presently proved. The conjecture is a testable
research target, not a renamed conclusion.

\section{Barrier audit and parallel routes}

A complete VIOT proof must identify why its decisive step does not relativize,
why the useful property is not simultaneously large and efficiently
constructive in the natural-proofs sense, and why low-degree algebraic
extension does not preserve the argument. It must also distinguish Boolean
from arithmetic circuit models and prove superpolynomial, not merely linear,
growth.

The broader program runs the following routes in parallel:
\begin{enumerate}[leftmargin=*]
\item unrestricted gate elimination and higher rewriting;
\item algorithms-to-lower-bounds transfers \cite{Williams};
\item constructive refuters and range avoidance;
\item proof complexity and bounded arithmetic, including recent consistency
results for EXP lower bounds \cite{AM};
\item algebraic geometry and representation theory;
\item equivariant topology and VIOT;
\item an adversarial search for a worst-case polynomial-time SAT algorithm;
\item meta-complexity and minimum circuit size.
\end{enumerate}
A route is retired only when a counterexample or incompatibility theorem
addresses its exact conjecture. The absence of a counterexample is not a
proof.

\section{Lean formalization}

The companion Lean project formalizes the face-closure counterexample, the
absent many-one reduction, nonadditivity of \(\rho\), failure of the proposed
contraction, verifier-padding equivalence, abstract obstruction soundness, and
abstract class-separation criteria. CI rejects `sorry', `admit', and
user-declared `axiom'. These are formalized auxiliary results; they are not a
formalization of \(P\ne NP\).

\section{Conclusion}

The audited homological construction fails for explicit minimal reasons and
does not resolve \(P\) versus \(NP\). The failure exposes a general design
constraint: a complexity-topological invariant must survive
resource-preserving changes of computational presentation.

VIOT turns that constraint into a localization-first program and records the
exact circuit-soundness and target-nonvanishing bridges still needed. The next
credible milestone is a calibrated theorem recovering a known restricted
lower bound from a presentation-invariant obstruction. Only after that
calibration should the program target NP versus \(P/poly\).

\begin{thebibliography}{99}
\bibitem{Cook} S. Cook, The complexity of theorem-proving procedures, STOC, 1971.
\bibitem{BGS} T. Baker, J. Gill, R. Solovay, Relativizations of the \(P=?NP\) question, \emph{SIAM J. Comput.} 4 (1975).
\bibitem{RR} A. Razborov, S. Rudich, Natural proofs, \emph{JCSS} 55 (1997).
\bibitem{AW} S. Aaronson, A. Wigderson, Algebrization: a new barrier in complexity theory, \emph{ACM TOCT} 1 (2009).
\bibitem{Williams} R. Williams, Nonuniform ACC circuit lower bounds, \emph{JACM} 61 (2014).
\bibitem{CDJ} M. Carmosino, N. Dang, T. Jackman, Convergent gate elimination and constructive circuit lower bounds, arXiv:2602.17942, 2026.
\bibitem{RenWilliams} H. Ren, R. Williams, Near-maximum circuit lower bounds for exponential time with Merlin--Arthur queries, arXiv:2607.09963, 2026.
\bibitem{AM} A. Atserias, M. M\"uller, From G\"odel incompleteness to the consistency of circuit lower bounds, arXiv:2604.25251, 2026.
\bibitem{FHV} F. Frick, K. Hosseini, A. Vasileuski, A \(\mathbb Z_2\)-topological framework for sign-rank lower bounds, arXiv:2604.01510, 2026.
\bibitem{Tang} J.-G. Tang, A homological separation of \(P\) from \(NP\) via computational topology and category theory, arXiv:2510.17829v2.
\end{thebibliography}

\end{document}
